\documentclass[11pt]{article}
\usepackage[margin=2.4cm]{geometry}
\usepackage{amsmath,amssymb,booktabs,hyperref,xcolor,listings,enumitem}
\usepackage[T1]{fontenc}
\usepackage{lmodern}
\hypersetup{colorlinks=true,linkcolor=blue!50!black,urlcolor=blue!50!black,citecolor=blue!50!black}
\lstset{basicstyle=\ttfamily\small,columns=fullflexible,frame=single,framerule=0.3pt,backgroundcolor=\color{gray!6},
  keywordstyle=\color{blue!60!black},commentstyle=\color{gray!70!black},language=Python,breaklines=true}
\setlist{itemsep=2pt,topsep=3pt}
\newcommand{\vla}{\pi_\theta}

\title{Code-as-Steering:\\ Continual Programmatic Adaptation of a Frozen Vision-Language-Action Model}
\author{ASPIRE project --- working document}
\date{\today}

\begin{document}
\maketitle

\begin{abstract}
A coding agent builds a library of executable programs that change how a \emph{frozen} vision-language-action model (VLA)
behaves at inference time. A program reads the observation and, crucially, the VLA's own internal state (activations,
attention), and decides which internal interventions to apply: additive steering directions in chosen layers, ablations,
attention edits, or, as a last resort, small adapters. The agent discovers these interventions by controlled experiments in
simulation, verifies them on held-out situations, and stores the verified ones with their applicability conditions so they
can be reused and composed on new tasks. The VLA weights never change, so every gain is attributable to the library. This
document states the idea, the implementation on SmolVLA in LIBERO and LIBERO-PRO, the first measurements and the plan.
\end{abstract}

\section{Idea}

Two lines of work meet here. \emph{Code-as-policies} systems (including our own ASPIRE) let a language model write the
robot's controller; they adapt fast but replace the learned policy with hand-written perception and motion code.
\emph{Activation steering} shows that adding a direction to a network's hidden states can change its behaviour without
training; for VLAs this has been demonstrated with fixed, hand-picked directions. What is missing is a system that
\emph{discovers} which interventions help, \emph{when} to apply them, and keeps that knowledge as reusable code.

\paragraph{Question.} Can a coding agent autonomously discover, verify and compose programs that control the internal
behaviour of a frozen VLA, so that the robot acquires new capabilities continually, without fine-tuning?

\paragraph{Why it could work.} A fine-tuned VLA already contains most of the skills a perturbed task needs (grasp, place,
open); failures under distribution shift (moved objects, paraphrased instructions) often look like the right skill
applied to the wrong target, or a commitment that never happens. Such failures may be correctable by small changes of
internal state that a program can trigger at the right moment, and the simulator gives cheap, exact feedback on whether
a change helped (an episode takes 2--10\,s here).

\section{Formulation}

Let $\vla$ be the frozen VLA. At each replanning step it maps the observation $o_t$ (two images, proprioception) and the
instruction $\ell$ to an action chunk $a_{t:t+H}$. Write $h_{l}$ for the hidden state that layer $l$ adds to the residual
stream. Three families of interventions, from cheapest to most expensive:
\begin{align}
\text{activation:}\quad & h_l' = h_l + \alpha(s)\,v_l, \qquad \|v_l\|=1, \label{eq:act}\\
\text{parameter:}\quad & W' = W + \textstyle\sum_i g_i(s)\,B_iA_i \quad\text{(low-rank adapters)}, \\
\text{action:}\quad & a_t' = a_t + \Delta a(s_t).
\end{align}
Here $s$ is everything a program may read: the observation, the episode so far, and \emph{the VLA's internals}
(Section~\ref{sec:introspect}). The first version keeps $\theta$ frozen and uses only~\eqref{eq:act} and simple
action edits, so every effect is attributable to the library.

\paragraph{Library entry.} An entry is executable knowledge $\mathcal S_i=(p_i,f_i,q_i)$: an applicability condition
$p_i(s)$, a steering program $f_i(s)$ that returns interventions, and evidence $q_i$ (success with and without the entry, on
which tasks and initial states, at what simulator cost). A program can call other entries, so entries compose.

\paragraph{Closed-loop internal steering.} Because $p_i$ and $f_i$ may read internal signals, a program can say for
example: \emph{if the action tokens attend to the wrist image less than usual and the projection on the ``grasp
committed'' direction is low, push that direction in the middle VLM layers}. The agent writes such rules at the level of
detail of assembly code, and the simulator decides whether they help.

\section{Implementation}

\subsection{Policy and simulator}
SmolVLA (450M parameters, checkpoint \texttt{lerobot/smolvla\_libero}, fine-tuned on the four LIBERO suites). It has two
stacks: 16 layers of a SmolVLM2 text model (width 960) that encode a \emph{prefix} of agentview-image, wrist-image,
instruction and state tokens once per chunk, and 16 layers of an \emph{action expert} (width 720) that denoise a chunk of
$H=50$ actions in 10 flow-matching steps, attending to the prefix (cross-attention, with self-attention every second
layer). The rollout copies lerobot's LIBERO evaluation exactly: $256{\times}256$ images rotated by $180^\circ$, state
$=$ end-effector position, axis-angle and gripper (8 numbers), initial state $e \bmod 50$ for episode $e$, 10 settling steps,
the suite's step limit. Environments come from LIBERO-PRO, so the same code runs the trained suites and their perturbed
versions (object positions swapped, paraphrased instructions, new objects).

\subsection{Intervention sites}
SmolVLA runs its layers in a custom loop, so whole-layer hooks never fire; each layer's MLP is still called as a module and
its output is added to the residual stream. A forward hook on \texttt{layer.mlp} therefore implements~\eqref{eq:act}
exactly. VLM-layer interventions act once per chunk on the prefix (and through the cached keys and values on every
denoising step); expert-layer interventions act on every denoising step. Strengths are expressed in units of the spread of
activations along the direction (``a $k\sigma$ push'').

\subsection{Introspection}\label{sec:introspect}
Per chunk, a small summary that programs read and traces store:
\begin{itemize}
\item \textbf{attention}: for each expert layer, the share of the action tokens' attention on each prefix segment (agentview,
wrist, instruction, state) and on the action tokens, averaged over heads and denoising steps;
\item \textbf{mlp norms}: norm of each layer's mean MLP output;
\item \textbf{projections}: projection of each layer's mean MLP output on named directions from the library;
\item \textbf{behaviour}: end-effector trajectory, gripper and executed actions (every fifth step).
\end{itemize}

\subsection{Directions}
Directions are contrasts of recorded activations (mean over tokens, per layer):
\emph{success} --- successful minus failed episodes of a task; \emph{shift} --- a trained task minus its perturbed version
(``pull the representation back towards the training distribution''); later, instruction pairs and directions found by
the agent itself.

\subsection{Offline lab}
Simulator rollouts are the ground truth but cost seconds; most questions about the model can be asked offline.
\texttt{scripts/run\_program.py --save-obs} stores the raw observation at every chunk start and every executed action. The lab
(\texttt{steer/lab.py}) then predicts the 50-action chunk for any saved observation in about 0.5\,s with a \emph{fixed}
flow-matching noise seed, so any change of the chunk is caused by the intervention. It captures per-token activations of any
layer, edits them, and \emph{replaces} them with activations captured on another observation or another instruction (activation
patching and swapping, by token span: agentview, wrist, instruction, state).

\subsection{Neurons}
The MLP neurons are the inputs of each layer's \texttt{down\_proj} (2560 per VLM layer). Two tools:
\begin{itemize}
\item \textbf{Attribution patching}: for any differentiable function $f$ of the (normalised) action chunk, the first-order effect
of neuron $n$ is $\sum_{\text{tokens}} a_n\,\partial f/\partial a_n$, computed in one backward pass through the ten denoising
steps for all neurons of the chosen layers.
\item \textbf{Dosed pushes}, with the recipe used for LLM neuron steering: the added activation is
\[
\Delta a_n = \text{dose}\times\frac{\operatorname{median}_{\text{tokens}\,1:}\,\|r_L\|}{\|W_{\text{down}}[:,n]\|},
\]
where $r_L$ is the residual stream leaving layer $L$, so a dose of one moves the residual by about one typical norm along the
neuron's output direction.
\end{itemize}

\subsection{Classic interpretability}
\texttt{steer/interp.py} builds an activation bank over saved episodes: for every chunk, the activations of all neurons of
chosen layers, summarised per prefix segment (VLM) or over action tokens and denoising steps (expert), with labels (success,
gripper phase, end-effector position, executed displacement and gripper command, time). On it: neuron selectivity (AUROC for binary
labels, correlation otherwise), max-activating examples, word-selective neurons on the instruction tokens, and linear probes with
folds grouped by episode, whose weight vectors double as steering directions. A first check on 74 chunks: gripper phase is linearly
decodable from wrist-image tokens at VLM layer 8 (accuracy 0.81 against 0.69 chance) and the upcoming vertical motion from the state
token at layer 12 ($R^2=0.72$); one neuron of layer 8 responds strongly and specifically to the word ``cheese''.

\subsection{Adapters}
\texttt{steer/lora.py} wraps chosen linear layers of either stack with zero-initialised low-rank adapters, trains them with
SmolVLA's own flow-matching loss on saved (observation, executed chunk) pairs (40 steps take about 3\,s), and detaches them so
the base model is restored exactly. Programs may load an adapter in \texttt{setup} and remove it in \texttt{on\_end}.

\subsection{Programs and library}
Programs are deliberately \emph{open}: the agent can log anything and intervene with any inference-time method. A program
file may define four optional functions, and the context exposes the whole model:
\begin{itemize}
\item \texttt{setup(ctx)} once per episode: register hooks on any module (\texttt{ctx.hook}), patch methods such as the
attention function (\texttt{ctx.patch}), rewrite the instruction (\texttt{ctx.instruction});
\item \texttt{program(t, obs, internals, ctx)} before every action chunk: return additive steering terms;
\item \texttt{on\_step(t, obs, action, ctx)} every control step: replace the action;
\item \texttt{on\_end(result, ctx)}: undo in-place edits.
\end{itemize}
\texttt{ctx.vla}, \texttt{ctx.vlm\_layers} and \texttt{ctx.expert\_layers} give the model, \texttt{ctx.log(key, value)} writes
arbitrary values into the episode trace, \texttt{ctx.memory} keeps state within the episode. The only hard constraints:
no simulator ground truth (object poses, simulator state), and the frozen weights must be identical after each episode
(checked by the evaluator; adapters are allowed as separate modules loaded and removed by the program).
\begin{lstlisting}
DIRECTIONS = ["goal_lan_0_shift_vlm"]       # projected into `internals`

def setup(ctx):                             # log per-token norms of VLM layer 8 on every prefix pass
    ctx.hook(ctx.vlm_layers[8].mlp,
             lambda m, i, out: ctx.log("vlm8_norms", out[0, :8].norm(dim=-1)))

def program(t, obs, internals, ctx):        # pull the paraphrase towards the trained instruction until the grasp
    if obs["robot0_gripper_qpos"][0] < 0.03:
        return []
    return [ctx.push("goal_lan_0_shift_vlm", layer=7, sigmas=2.0)]

def on_step(t, obs, action, ctx):           # slow down once something is held
    if obs["robot0_gripper_qpos"][0] < 0.03:
        action = action.copy(); action[:3] *= 0.5
        return action
\end{lstlisting}
\texttt{scripts/run\_program.py} evaluates a program on development initial states (episodes 0--19) and writes traces;
held-out initial states (20--49) are only used by \texttt{scripts/verify.py}, which promotes a program into the library when it
beats the frozen VLA there (one-sided Fisher exact test, 5\%).

\subsection{Agent loop}
Two modes share the same tools and library.
\paragraph{Research agent (main mode).} A Claude session runs unattended for hours (\texttt{prompts/explorer.md},
\texttt{scripts/explore.py}). Its mission is to make the frozen VLA solve perturbed tasks by dissecting how the activations work,
running causal experiments (patching, swapping, ablations, neuron pushes, toy probes with synthetic images or text-only prefixes,
toy adapters), turning findings into programs, combining and stitching them, and keeping what works as code in
\texttt{library/steerlib/} with its evidence. It keeps a lab notebook (\texttt{library/notes/JOURNAL.md}) from which a restarted
session continues. Ideas are screened on 5--10 development episodes (initial states 0--19); a promising program is confirmed on 20
and submitted; the driver verifies it on held-out initial states 20--49 against the frozen VLA (one-sided Fisher exact test, 5\%)
and records the result in \texttt{library/INDEX.md}.

\paragraph{ASPIRE-style round.} For one target task (\texttt{prompts/agent.md}): three candidate programs built from the library,
up to six repairs from their traces, then a library update and held-out verification.

\paragraph{Rules.} Development episodes only; no simulator ground truth inside programs and no benchmark task name; for
paraphrased-instruction tasks, no use of the task's original wording (an instruction rewrite must be a general rule); frozen
weights identical after every episode. Escalation is free: the agent decides when logging, activation and attention edits,
neuron pushes, action edits or adapters are worth it, and every entry records its simulator cost so methods can be compared at
equal budget. A verified program can later be \emph{compiled} into an adapter by distilling its behaviour.

\section{First measurements}\label{sec:results}
\begin{table}[h]\centering\small
\setlength{\tabcolsep}{4pt}
\begin{tabular}{lccccccccccc}
\toprule
suite & 0 & 1 & 2 & 3 & 4 & 5 & 6 & 7 & 8 & 9 & mean \\
\midrule
\texttt{libero\_goal} & 1.0 & 1.0 & 1.0 & 0.5 & 0.9 & 0.6 & 0.4 & 1.0 & 1.0 & 1.0 & 0.84 \\
\texttt{libero\_goal\_lan} & 0.2 & 0.7 & 0.7 & 0.0 & 0.8 & 0.0 & 0.2 & 0.1 & 0.1 & 0.5 & 0.33 \\
\texttt{libero\_goal\_swap} & 0.0 & 0.0 & 0.0 & 0.0 & 0.0 & 0.0 & 0.0 & 0.0 & 0.0 & 0.0 & 0.00 \\
\texttt{libero\_spatial} & 1.0 &  &  &  &  &  &  &  &  &  & 1.00 \\
\bottomrule
\end{tabular}
\caption{Frozen SmolVLA (\texttt{lerobot/smolvla\_libero}): success rate per task (10 episodes, initial states 0--9). \texttt{\_swap}: object positions swapped; \texttt{\_lan}: paraphrased instruction.}
\end{table}

No steering sweep yet.


\paragraph{The policy is sensitive to wording and memorises positions.} On the trained suite the frozen policy reaches 0.84; with
paraphrased instructions success ranges from 0.0 to 0.8 per task; with object positions swapped it is 0.0 on every task tried so
far, i.e.\ the policy goes to the trained positions.

\paragraph{Pilot: plain mean-difference steering has no effect.} The direction ``trained minus paraphrased'' (mean MLP output over
all prefix tokens and chunks) added at VLM layers 3, 7, 11, 15 and expert layer 3 with strengths up to $4\sigma$ left success on
\texttt{libero\_goal\_lan} task 7 at 0.00--0.05 (held-out states, 20 episodes per setting). The literature review gives the likely
reasons: the direction is averaged over 128 image tokens that are identical in both conditions, the strength unit (spread of
episode means) is tiny relative to per-token norms, and it is added to every token. Mechanism tests (direct steering of large random
directions) confirm the hooks do change behaviour: pushes of 500 along a random direction break a task the policy normally solves.

\section{Results of the first research session}\label{sec:session1}
One unattended session of the research agent (2026-10-26, 16:40--00:40, Claude Opus, one GPU) on the paraphrased and swapped
goal and spatial suites. Every claim below is either a held-out verification (driver, initial states 20--49) or an experiment
recorded in the agent's notebook (\texttt{library/notes/JOURNAL.md}, experiment numbers E$n$) and summarised in
\texttt{library/knowledge/network\_map.md}.

\begin{table}[h]\centering\small
\begin{tabular}{llllll}
\toprule
task & program & kind & frozen & program & $p$ \\
\midrule
\texttt{libero\_goal\_lan 0} & \texttt{canon\_v5} & text rewrite & 11/30 & 26/30 & 0.000 \\
\texttt{libero\_goal\_lan 1} & \texttt{lexfacts} & in-network edit (stored fact) & 27/30 & 30/30 (n.s.) & 0.119 \\
\texttt{libero\_goal\_lan 2} & \texttt{lexfacts} & in-network edit (stored fact) & 22/30 & 26/30 (n.s.) & 0.167 \\
\texttt{libero\_goal\_lan 6} & \texttt{canon\_v5} & text rewrite & 6/30 & 18/30 & 0.002 \\
\texttt{libero\_goal\_lan 7} & \texttt{lexedit\_live} & in-network edit (live) & 1/30 & 29/30 & 0.000 \\
\texttt{libero\_goal\_lan 8} & \texttt{lexfacts\_mask} & in-network edit (stored fact) & 0/30 & 25/30 & 0.000 \\
\texttt{libero\_goal\_lan 9} & \texttt{lexedit\_live} & in-network edit (live) & 15/30 & 26/30 & 0.002 \\
\texttt{libero\_goal\_swap 9} & \texttt{approach\_handback} & hybrid hand-back & 0/30 & 4/30 (n.s.) & 0.056 \\
\texttt{libero\_spatial\_lan 0} & \texttt{canon\_v7} & text rewrite & 10/30 & 15/30 (n.s.) & 0.147 \\
\texttt{libero\_spatial\_lan 1} & \texttt{canon\_v6} & text rewrite & 0/30 & 22/30 & 0.000 \\
\texttt{libero\_spatial\_lan 2} & \texttt{canon\_v7} & text rewrite & 3/30 & 10/30 & 0.029 \\
\texttt{libero\_spatial\_lan 3} & \texttt{canon\_v7} & text rewrite & 1/30 & 24/30 & 0.000 \\
\texttt{libero\_spatial\_lan 4} & \texttt{canon\_v13} & text rewrite & 0/30 & 21/30 & 0.000 \\
\texttt{libero\_spatial\_lan 5} & \texttt{canon\_v6} & text rewrite & 0/30 & 24/30 & 0.000 \\
\texttt{libero\_spatial\_lan 6} & \texttt{canon\_v6} & text rewrite & 0/30 & 27/30 & 0.000 \\
\texttt{libero\_spatial\_lan 7} & \texttt{canon\_v12} & text rewrite & 0/30 & 24/30 & 0.000 \\
\texttt{libero\_spatial\_lan 8} & \texttt{canon\_v13} & text rewrite & 0/30 & 19/30 & 0.000 \\
\texttt{libero\_spatial\_lan 9} & \texttt{canon\_v12} & text rewrite & 0/30 & 20/30 & 0.000 \\
\texttt{libero\_spatial\_swap 0} & \texttt{carry\_handback4} & hybrid hand-back & 0/30 & 8/30 & 0.002 \\
\texttt{libero\_spatial\_swap 2} & \texttt{carry\_handback2} & hybrid hand-back & 0/30 & 20/30 & 0.000 \\
\texttt{libero\_spatial\_swap 4} & \texttt{carry\_handback2} & hybrid hand-back & 0/30 & 11/30 & 0.000 \\
\texttt{libero\_spatial\_swap 5} & \texttt{carry\_handback2} & hybrid hand-back & 0/30 & 8/30 & 0.002 \\
\texttt{libero\_spatial\_swap 7} & \texttt{carry\_handback4} & hybrid hand-back & 0/30 & 18/30 & 0.000 \\
\texttt{libero\_spatial\_swap 9} & \texttt{carry\_handback4} & hybrid hand-back & 0/30 & 17/30 & 0.000 \\
\bottomrule
\end{tabular}
\caption{Held-out verification (initial states 20--49, 30 episodes): best program per task against the frozen policy on the same states, one-sided Fisher exact test. 20 of 24 tasks promoted.}
\end{table}


\paragraph{Verified gains.} 21 tasks were promoted: 15 paraphrased-instruction tasks (\texttt{goal\_lan} 0, 6, 7, 8, 9;
\texttt{spatial\_lan} 1--9) and, for the first time, 6 swapped-position tasks (\texttt{spatial\_swap} 0, 2, 4, 5, 7, 9), where the
frozen policy scores 0/30 everywhere. \texttt{goal\_swap} is still essentially unsolved (4/30 on task 9, not significant), and
\texttt{goal\_lan} 3 and 5 are unsolved (they need a word or a two-step skill that only the target task's own training taught).

\paragraph{Three kinds of programs.}
\begin{itemize}
\item \textbf{Text rewrite} (\texttt{steerlib/canon.py}): a noisy-channel translation of the instruction into the policy's training
vocabulary, with the base SmolVLM2 (and Llama-2-7B) as a meaning check and the target task's twin excluded from all statistics. On
eight fresh paraphrases written before any result it raised mean success from 0.25 to 0.73.
\item \textbf{In-network edits, text unchanged} (\texttt{steer/edit.py}): a closed-form rank-one edit of the VLM layer-1 down
projection that makes the layer write the value of a known verb for an unknown one (``switch'' $\to$ ``turn''), fitted live on the
episode's first observation: \texttt{goal\_lan} 7 from 1/30 to 29/30 and \texttt{goal\_lan} 9 from 15/30 to 26/30, equal to the text
rewrite. A ``place $\to$ put'' fact learned on one sentence transfers to other sentences and tasks (\texttt{goal\_lan} 8 from 0/30 to
25/30, with the deletion of out-of-vocabulary words through the padding mask).
\item \textbf{Hybrid hand-back} for swapped layouts: the VLA grasps; the program carries the object above the destination tracked
in the image with the policy's own SigLIP patch features (voting over reference frames), then hands control back with a
virtual-frame state so the VLA's own final alignment and release run; a nudge-and-retry search when it hovers. This is closer to
code-as-policies than to steering, and is reported as such.
\end{itemize}

\paragraph{What the network does (mechanistic findings).}
\begin{itemize}
\item \textbf{Paraphrase failures are lexical.} Unknown verbs are represented like object nouns (``switch'' near ``compartment'',
``drawer'') and the possessive like the end of the instruction; known words stay at cosine 0.93--1.00 with their trained context at
every layer. Fine-tuning removed general English from the 16-layer VLM (E02).
\item \textbf{How the action expert reads the instruction:} as (token keys and values, rotary position) pairs through its
cross-attention (odd) layers; the verb lives in its own token's KV at every VLM layer, there is no summary token. Patching only the
verb's residual restores 96\% of a rewrite's effect after layer 0, 50\% after layer 8 and none after layer 15 (E15, E22, E23).
\item \textbf{Hence where to edit:} the layer-1 down projection, where a closed-form edit fitted on two observations reaches the
patching bound on new observations (E16). Edits can substitute a word's meaning but cannot insert or reorder words; structural
paraphrases still need the text rewrite.
\item \textbf{Word identity is dense} (more than 256 neurons to reconstruct one correction, E19), and function words matter as much
as verbs: wordings are memorised as exact patterns (E25).
\item \textbf{Targets are memorised per instruction.} The images steer \emph{how} to move (blinding the wrist camera changes the path
tenfold) but not \emph{where}; probe directions, state shifts and attention boosts towards the images do not move the target (E06,
E10--E12, E17). In goal scenes the destination is a smooth function of the destination noun's representation: interpolating
``plate'' towards ``stove'' at that token moves the endpoint along the line between them (E20); in spatial scenes it is hard-wired
(E21).
\item \textbf{A single action-expert neuron controls gripper timing} (expert layer 4, neuron 1855): dose $-1$ keeps the gripper open,
$+2$ closes it at once, $+1$ does nothing (threshold behaviour, E28).
\end{itemize}

\paragraph{Learning and sample efficiency.}
\begin{center}\small
\begin{tabular}{p{3.6cm}p{3.6cm}p{3.6cm}p{3.6cm}}
\toprule
component & learns & data & result \\
\midrule
text rewrite & translation into the training vocabulary & none (a language model and the other trained instructions) & 9 spatial and 3 goal tasks promoted \\
stored lexical fact & unknown word = known word at VLM layer 1 & one sentence pair, two observations & transfers to other sentences and tasks \\
live lexical edit & the same, fitted at run time & the episode's first observation & \texttt{goal\_lan} 7: 29/30, 9: 26/30 \\
LoRA from own success & paraphrase $\to$ behaviour & one self-generated success, about 15\,s of training & \texttt{spatial\_lan} 1: 3/5 (frozen 0/10); no transfer across tasks \\
LoRA for a swapped layout & new layout & seven self-generated hybrid successes & the VLA alone 1/10 (program 8/15): not enough \\
\bottomrule
\end{tabular}
\end{center}

\paragraph{Caveats.} (i) Several rewrites coincide with the original trained wording of the target (``turn on the stove''); every
word comes from the paraphrase or from \emph{other} trained sentences, but this makes the language results easier than they look.
(ii) Two choices are pending review: the fixed agentview camera model from the LIBERO scene definition is used as a sensor
calibration constant (no object poses), and trained-suite recordings of the same task serve as references for swapped layouts.
(iii) Before \texttt{libero\_object\_lan} was reserved, the agent had seen its instructions; the transfer comparison was moved to
the untouched \texttt{libero\_10\_lan}. (iv) One session, one checkpoint; development episodes were screened many times, which is why
only the held-out numbers are claimed.

\section{Plan}
\begin{enumerate}
\item \textbf{Done:} baseline; tools (lab, neurons, interpretability, edits, adapters, open programs); first research session
(Section~\ref{sec:session1}).
\item \textbf{Library transfer} (running next): the same agent and budget (2\,h) on two unseen \texttt{libero\_10\_lan} tasks, once
with the seed library and once with the grown library; repeat on more unseen tasks if the signal is clear.
\item \textbf{Equal-budget learning baseline}: plain LoRA fine-tuning with the same number of simulator episodes and successes that a
library fix consumed, to measure sample efficiency against standard training.
\item \textbf{In-network structural fixes}: insert tokens into padding slots with assigned rotary positions (the expert reads
(token, position) pairs), so insertions and reorderings can be fixed inside the network, not only substitutions.
\item \textbf{Swapped layouts}: perception (re-tracking, every object's template), wrist-camera alignment before the hand-back,
internal levers for the target (the destination-noun interpolation found in goal scenes), and adapters with more data.
\item \textbf{Compilation}: distil verified programs and edits into adapters, and retrieve adapters by activation (adapter bank).
\end{enumerate}

\paragraph{Baselines and ablations.} Frozen VLA; fixed steering (no state dependence); state-dependent steering without
internal signals; programs without the library (search per task); full library. An equal-budget comparison against learned
residual corrections or adapters trained with the same number of simulator episodes is needed to show that the benefit comes
from program synthesis and not from simulator experience alone.

\paragraph{Risks.} (i) Steering directions may not be causal or may not exist for a failure mode; the sweep in step~2 tests
this before any agent work. (ii) Gains on development initial states may be noise; verification on held-out states and
binomial confidence intervals guard against it. (iii) The perturbed suites may be failures of perception that no internal
push can fix; language perturbations are the most likely to respond to steering, position swaps the least.

\section{Related work}
A survey (\texttt{library/LITERATURE.md}) and an independent review against the sources (\texttt{library/LITERATURE\_REVIEW.md})
seed the library; the points that matter most here:
\begin{itemize}
\item \textbf{VLA internals.} Clamping clusters of feed-forward neurons in OpenVLA and $\pi_0$ changes motion style (speed,
height) but no task-success gains are reported (Häon et al., CoRL 2025). Optimal-transport maps on the second-to-last action-expert
layer, tested on \texttt{smolvla\_libero}, control style reliably while mean-difference steering was inconsistent (DiMaS).
Conceptor gates on action-expert residuals from about 15 success/failure rollouts per task raise success on LIBERO-10, measured on
deliberately early, weak checkpoints (COAST). Activations injected into the action expert replay trajectories in absolute
coordinates, while goal selection lives in the VLM: change \emph{where} through the prefix or the sampler and \emph{how} through
the expert.
\item \textbf{Sampler guidance.} Reward-gradient guidance of the flow velocity with particle resampling is the only method with
LIBERO-PRO results ($\pi_{0.5}$ 23.7\% to 36.8\%, VLS); its keypoints come from segmentation and depth, so for us the reward must
come from the camera images and proprioception. Action coherence guidance (identity attention in some action-expert self-attention
layers, then extrapolation) has been shown on SmolVLA hardware; in our model the self-attention layers are expert layers
$0,2,\dots,14$. Note LeRobot's convention: $x_t = t\,\epsilon + (1-t)\,a$ with $t$ from 1 (noise) to 0, so the clean estimate is
$\hat a = x_t - t\,v_t$.
\item \textbf{Gating and verification.} Steering only when a failure detector fires beats always-on steering ($\pi_{0.5}$ 42.7 vs
33.8 with a learned latent policy, RL$^2$-VLA, which also needs a verifier for best-of-$N$).
\item \textbf{Attention.} Re-weighting attention towards instruction tokens (IGAR, and PASTA for LLMs) improves language grounding
at some cost in success under contradictory instructions.
\item \textbf{LLM steering.} Contrastive activation addition, inference-time intervention (mass-mean shift), representation
engineering, directional ablation and conceptors are the building blocks; benchmarks find prompting and low-rank representation
fine-tuning stronger than mean-difference and SAE steering, and steering vectors fragile under rewording.
\end{itemize}
The distinctive claim here is the automated discovery, verification, composition and reuse of state- and internal-signal-dependent
steering programs by a coding agent, with the frozen weights untouched; a full novelty check is still needed.

\end{document}
